Les problèmes émergents des domaines physiques, quantiques ou biologiques
deviennent de plus en plus coûteux que ce soit en resources et en temps. Le
calcul haute-performance peut résoudre les problèmes de temps via la
paralléllisation ou l'utilisation d'architecture moderne (GPUs, accélérateur).
Pour l'aspect consommation de mémoire, les expressions d'algébre linéaire
peuvent simplifier, voir réduire la taille des objets traités, de telle sorte
que le problème soit calculable sur un ordinateur et correct algébriquement.
Ces méthodes sont très exploitées dans le domaine matriciel mais très peu dans
le domaine tensoriel.

Le projet de recherche proposé est d'étendre l'aspect haute-performance aux
applications tensorielles. Afin de réduire leur empreinte mémoire par les
méthodes d'algébre matricielle et/ou tensorielle. Cela reviendra à utiliser des
méthodes d'approximation et de compression, dites stables, sur différents types
de tenseur, que ce soit en précision uniforme ou mixte.

Une autre approche sur la façon de réduire l'utilisation de la mémoire peut
être envisagé pour les matrices creuses. Cette idée consiste à stocker
uniquement les indices de lignes ou de colonnes mais pas les deux à la fois.
Ainsi un vecteur d'index sera économisé. L'intuition est d'utiliser un vecteur
de vecteurs pour stocker les indices où chaque vecteur représentera soit les
lignes ou les colonnes. Ainsi une unicité des valeurs est maintenue dans chaque
vecteur et surtout l'ordre des valeurs n'importe plus (rendant beaucoup plus
d'algorithme plus applicable).

\subsubsection*{Résumé :}

Mon projet de recherche propose d'améliorer les applications tensorielles en
utilisant des techniques de calcul haute-performance et de précision mixte.
\begin{itemize}
    \item[\textbf{Axe 1 :}] Étudier les calculs parallèles dans le domaine tensoriel, principalement via des approximations de rang faibles (HOSVD, etc...) stables et de l'utilisation de précision mixte (FP16-FP32, etc...).
    \item[\textbf{Axe 2 :}] Utilisation d'une nouvelle approche de stockage des éléments pour une matrice creuse, via un seul vecteur d'indices.
\end{itemize}

\subsection{Compression et précision mixte sur les tenseurs}

Les tenseurs~\cite{koba09} sont la généralisation des matrices à l'ordre
supérieur de deux, et sont très utilisés dans les domaines de
physiques~\cite{apda81, vate07}, quantiques, et dans le signal
processing~\cite{sdfh17, co14}. Le problème émanant des tenseurs vient de la
taille des éléments à stocker qui est exponnentielle par rapport à l'ordre du
tenseur, ce problème est appelé \textit{curse of dimensionality}. La
``malédiction de la dimensionalité'' peut être résolue en appliquant des
méthodes de compressions qui transforment le tenseur original en tenseurs de
petites dimensions pour être calculé sur ordinateur; comme le fait le format
CP~\cite{koba09}.

\subsubsection*{Approche sur les méthodes itératives}

La méthode pour obtenir un tenseur au format CP, nommé CP-ALS, est une méthode
itérative donc nous n'avons pas de garantie de convergence. Cette analyse de
convergence est encore à ce jour très étudiée~\cite{ylc25}, mais peu dans le
domaine de la précision mixte~\cite{yashzh22}. Une approche serait de faire les
premières étapes de la fonction CP-ALS en faible précision (e.g. \textit{half}
précision -- fp16) pour récupérer les premières valeurs singulières pour
comprimer le tenseur. Une fois l'itération courante au niveau de la tolérance
prescrite par la précision de calcul (e.g. $\epsilon \simeq 1e-5$ pour
\textit{half}), les prochaines itérations sont calculées en précision plus
haute et ainsi de suite jusqu'à atteindre la précision finale attendue.
L'analyse de convergence est l'assurance que la tolérance sera atteinte et
démontrée dans ces papiers~\cite{golub65,blmv25}. Il est intéressant de noter
que même si la convergence en faible précision demande plus d'itération,
puisque celles-ci sont en faible précision, elles sont moins coûteuses, que ce
soit en mémoire et en calcul, spécifiquement avec les nouvelles générations de
GPU tensor cores intégrant des calculs en précision mixte tels que le
FP16-FP32. Cette approche a été étudiée durant ma thèse mais sans aboutir par
manque de temps.

Une autre approche serait d'appliquer un sketch via les méthodes randomisées
sur le tenseur et potentiellement à chaque itération pour aller plus loin que
l'état de l'art~\cite{bbk18}.

\subsubsection*{Approche sur les méthodes directes}

D'autres formats de compression de tenseur existent tels que le
Tucker~\cite{tucker66, dedeva20b}, tensor-train~\cite{os11}, et hierarchical
Tucker~\cite{haku09, gras10} formats. Celles-ci peuvent s'obtenir depuis un
tenseur plein par une \textit{low-rank approximation}~\cite{koba09, os11}, tel
que HOSVD pour le Tucker format, TT-SVD pour tensor-train, et HT-SVD pour le
hierarchical Tucker. Chacune de ces méthodes sont directes, ce qui donne une
quasi-optimal approximation après un appel de ces fonctions. Ces méthodes
peuvent ainsi profiter d'un raffinnement itératif avec un contexte de précision
mixte commencant par de la faible précision jusqu'à atteindre la haute
précision. L'approximation restera stable par le fait que les approximations
sont révélateurs de rang, assurant ainsi une stabilité de ces méthodes.

Pour être plus précis, l'approche envisagée est de montrer que les méthodes
d'approximations sont équivalentes et spécifiques à la structure qu'ils ont en
entrée. Pour confirmer cette idée, nous pouvons nous concenter sur le format le
plus global, qui est le format Tucker hiérarchique (hierarchical Tucker), et
démontrer que sa méthode de compression HT-SVD est identique à celle du format
Tucker de base pour un Tucker hiérarchique avec un seul parent. Le principe est
le même pour le format tensor-train en partant du format Tucker hiérarchique.

Une fois cette idée confirmée et prouvée, nous peaufinnerons la méthode HT-SVD
avec des optimisations possibles tels que l'ordre des contractions des noeuds
du tensor. Ainsi nous pourrons étendre la méthode à des formats plus étendus
que les Tucker hiérarchique en comprimant certains de leur noeuds pour nous
retrouver dans un cas favorable.

\subsection{Nouveau framework sur les matrices creuses}

Les matrices creuses sont très utilées en informatique pour leur faible
consommation mémoire et les formats structurés des matrices qu'ils décrivent.
Actuellement, l'optimisation mémoire est de stocker uniquement les indices
lignes $i$ et colonnes $j$ tels que la matrice est non nul $A_{i,j} \neq 0$. La
contrainte de cette façon est de stocker deux vecteurs d'indices $v_i$, $v_j$
avec les indexes ordonnés selon les lignes ou colonnes et un même indice donné
pour ne pas avoir de corruption de lecture. C'est-à-dire, si $A_{i,j} \neq 0$
alors $i$ et $j$ doivent être au même indice $k$ dans $v_i$ et $v_j$
respectivement, sinon si $v_j[k] = j+1$ alors $A_{v_i[k],v_j[k]} = A_{i,j+1}$
qui peut être nulle. Le désavantage de cette technique est l'ajout d'un élément
dans la matrice, il faut insérer l'index de cet élément $l$ dans les deux
vecteurs à l'indice $k_l$; or l'insertion d'un élément dans un vecteur est très
coûteux.

L'approche proposée est d'utiliser un vecteur possédant lui même les vecteurs
correspondant aux indices de la dimension la plus grande. Ainsi pour une
matrice $A\in \mathbb{R}^{5\times 4}$ nous avons la configuration suivante : $$
    A =
    \begin{bmatrix}
        a_{0,0} & 0       & 0       & 0       \\
        0       & 0       & a_{1,2} & 0       \\
        0       & 0       & 0       & 0       \\
        0       & a_{3,1} & 0       & 0       \\
        0       & 0       & a_{4,2} & a_{4,3} \\
    \end{bmatrix}
    \Rightarrow
    v_{ji} = [[0],[3],[1,4],[4]] \text{ ou } v_{ij} = [[0],[2],[],[1],[2,3]]
    % \begin{cases}
    % v_i = [0,1,3,4,4]\\
    % v_j = [0,2,1,3,4]
    % \end{cases}
$$,
comparé au résultat précédent qui serait $v_i = [0,1,3,4,4], v_j = [0,2,1,2,3]$.

Les inconvénients de cette méthode sont les suivantes :
\begin{itemize}
    \item Le stockage reste dynamique, pas optimal en mémoire.
    \item Allocation inutile de vecteur vide, qui peut devenir un surcoût non négligeable
          sur de très grandes matrices.
\end{itemize}
À contrario, les avantages sont les suivants :
\begin{itemize}
    \item La taille de chaque sous-vecteur peut être préalloué au max d'éléments d'une
          dimension si nécessaire.
    \item L'ordre des éléments dans un sous-vecteur n'importe pas.
    \item L'ajout d'un élément non nul est beaucoup plus rapide et moins coûteux.
    \item Le format est plus propice à la parallélisation.
\end{itemize}

Pour montrer l'intuition derrière voici l'affichage du temps et de la mémoire
de l'exemple énoncé au-dessus. Nous pouvons constater que $v_i$ et $v_j$ ne
coûte que $192$ bytes au total comparé à $v_{ji}$ qui en coûte $368$. Mais pour
l'insertion d'un élément nous avons uniquement $96$ bytes comparé à $448$ pour
chaque vecteurs.

\begin{figure}
    \center
    \includegraphics[scale=0.5]{ExampleJulia.png}
\end{figure}

\subsection{Intégration dans l'équipe \equipe}

\begin{comment}
%% Strasbourg %%
Bien que le premier projet de recherche sur les calculs tensoriels ne s'intégre
pas nécessairement dans les axes principaux de l'équipe \equipe, le second
projet par contre se voit plus abordable. En effet, l'idée de revoir la façon
de stocker les éléments d'une matrice se rapproche très fortement des travaux
de Cédric Bastoul~\cite{glb20} sur la direction de lecture des éléments d'une
matrice.

Les aspects de parallélisation sur les deux projets peuvent être étudiés avec
Béranger Bramas, qui travaillant aussi sur le Projet Exasoft, possède les
compétences nécessaires sur les systèmes d'ordonnancements tels que StarPU pour
efficacement exécuter les tâches sur les tenseurs. Ayant principalement utilisé
StarPU, et participé aux AG Numpex, je pense que Béranger pourra élargir ma
vision sur les modèles générés à base de tâches en m'initiant à Specx le modèle
qu'il développe actuellement~\cite{cb25}.

Pour le projet sur les tenseurs, je pensais faire l'implémentation en C, mais
plus précisement en C++. En effet, les avantages de compilation et de nouvelles
options apportées par C++ me conforte dans une implémentation élégante et très
maniable du code. De plus, dans l'équipe il y a Jens Gustedt, un expert du
langage C voir même du C moderne, qui pourra m'aiguiller sur la façon dont je
voudrais faire les templates pour les tenseurs (voir même les contrats).

Au niveau de la vérification de code, mon domaine d'expertise n'est pas très
grand, c'est pourquoi je voudrais beaucoup collaborer avec Arthur Charguéraud
sur le sujet.
%% Strasbourg %%
\end{comment}

\begin{comment}
%% Paris %%
Durant ma thèse j'ai fréquemment travaillé avec l'équipe Pequan. J'ai participé
à des réunions d'équipe et ai pu voir arriver de nouveaux permanents tel que
Pierre Jolivet. Mes travaux de thèse sont en lien direct avec ceux de Théo Mary
sachant qu'il a été un de mes encadrants de thèse. La grande partie de mon
expérience sur la précision mixte et le calcul d'approximation provient de ses
travaux~\cite{bmp24,blmv25}. Cependant, les sujets étudiés durant ma thèse mais
non traités pourront être réalisés au sein de cette équipe. Avec la
collaboration de Fabienne Jézéquel et Stef Graillat, travaillant sur l'aspect
de validation numérique des applications et de leur robustesse aux faibles
précisions (avec \href{https://promise.lip6.fr/}{PROMISE}), je pourrais étendre
le projet de recherche sur les tenseurs sur des modèles de précision cible
(custom) avec une grande fiabilité numérique.

Le projet portant sur le stockage de matrice ne possède pas de lien direct avec
des membres de l'équipe mais reste très proche du domaine sachant qu'il vise à
optimiser le stockage mémoire et donc rendre aussi les transferts de données
moins coûteux, ce qui implique une meilleur efficacité pour des exécutions
parallèles. Pourtant, je pourrais tirer partie de leur application ou ``cas
test'' d'exécution pour vérifier si l'idée peut converger sur une nouvelle
méthode de stockage de matrice. Les expériences pourront être vues avec
Dominique Béréziat, travaillant dans l'assimilation de données et le traitement
d'images sur des applications météorologiques~\cite{gabc26} (avec l'outil
GloFM). Ses domaines d'études utilisent souvent des matrices creuses dont le
projet pourrait en tirer partie.

Pour le projet sur les tenseurs, je pensais faire l'implémentation en C, mais
plus précisement en C++. En effet, les avantages de compilation et de nouvelles
options apportées par C++ me conforte dans une implémentation élégante et très
maniable du code. De plus, Pierre Jolivet possède une grande connaissance dans
le langage C en tant que développeur principal de
\href{https://petsc.org/release/}{PETSc}, ainsi il pourra m'aiguiller sur la
façon de rendre le code portable pour tous les utilisateurs venant de
différents domaines d'application (physique, quantique, etc...).

Au sein de cette équipe je pourrais apporter une nouvelle vision de la
recherche à l'aide du calcul tensoriel. Ainsi, diverses collaborations avec des
applications physiques pourraient aboutir. Je pourrais aussi m'intégrer dans
les projets déjà existants tels que ceux présents dans le labo, à savoir
\href{https://promise.lip6.fr/}{PROMISE} ou
\href{https://petsc.org/release/}{PETSc}.
%% Paris %%
\end{comment}

% \begin{comment}
%% Lyon %%
Durant ma thèse, j'ai principalement enseigné pour des L3 et des M1/M2 dans des
unités d'enseignement telles que la programmation parallèle et distribuée, via
OpenMP et MPI. La partie théorique était souvent accompagnée d'une partie
pratique avec soit des applications physiques, comme la résolution d'EDP, soit
des architectures récentes telles les machines du cluster Jean Zay.

En discutant avec Michele Pagani, j'ai appris que l'équipe \equipe se tourne
vers des enseignements plus applicatifs. Or, durant mon post-doctorat j'ai créé
un projet pratique nommé ``Projet Calcul Scientifique : Matrices randomisées et
leur utilisations pour les approximations de rang faible'' en collaboration
avec un collègue de mon équipe, Ronan Guivarch. Grâce à ces expériences je
pense pouvoir aider l'équipe \equipe dans la création ou le soutien de projets
pratiques pour des licences 3 ou des masters, voire enseigner dans des unités
hors de mon domaine comme celles de la théorie des graphes ou de la dynamique
probabilistique des modèles discrets enseignée par Éric Thierry.

Ayant un bagage sur l'optimisation des performances, j'ai collaboré avec
Mathieu Faverge et Emmanuel Agullo pour améliorer les librairies Chameleon et
StarPU sur les matrices structurées. De plus, sur la résolution de matrice
creuse, j'ai aussi cotoyé des chercheurs tels que Patrick Amestoy, qui
travaille sur MUMPS. Ainsi, je pourrai fournir des nouveaux sujets d'études sur
des librairies open source et optimisées comme celle susmentionnée qui pourront
mener à des propositions de thèse ou de recherche sur des optimisations ou sur
des applications dans des domaines externes tels que la physique.

J'aimerais aussi favoriser la diversité des langages de programmation dans les
enseignements pour élargie les possibiltés de recherche des étudiants, qu'elles
soient locales ou à l'étranger. C'est pourquoi mon expertise dans les langages
comme C/C++, Julia, Matlab me conforte dans cette idée et se retranscrirait
bien dans le domaine de recherche de Michele Pagani.

Sur le plan de la recherche, j'aimerais beaucoup étendre le sujet du calcul
tensoriel et de ses approximations de rang faible au-delà de la programmation
parallèle en mémoire partagée, jusque dans la mémoire distribuée. C'est
pourquoi j'aimerais beaucoup collaborer avec Isabelle Guérin Lassous sur ce
projet. Je pense qu'on pourrait également aborder le projet sur le stockage des
matrices creuses en mémoire distribuée mais je n'ai pas encore creusé le sujet.

J'aimerais préciser que les deux projets ont pour objectif d'être parallélisés
et que, pour ce faire, je recommande d'utiliser un modèle \textit{task-based},
ce qui amènera à gérer un \textit{directed acyclic graph} (DAG)
comme dans les travaux de recherche de Loris Marchal.
%% Lyon %%
% \end{comment}

\printbibliography[keyword=pdr]